% SPDX-License-Identifier: Apache-2.0
%
% The project's dynamics: how the work actually went. Built by
% //docs:dynamics. One column, because the chart is wide.
\documentclass[journal,onecolumn,9pt]{IEEEtran}

\newcommand{\listingsize}{\normalsize}
% SPDX-License-Identifier: Apache-2.0
%
% The house preamble, shared by every document in this directory. Loaded
% after \documentclass, before \title. Keeping it in one file is what
% stops two articles drifting apart in font, listing style or figure
% style.
% Computer Modern. IEEEtran selects Times (ptm) by default, so the face has
% to be chosen explicitly.
%
% Latin Modern is the Computer Modern variety used here, and the choice is
% forced rather than aesthetic. Under T1 encoding, plain `cmr` has no Type 1
% outlines unless cm-super is installed, and it is not in the pinned
% distribution, so pdflatex falls back to EC bitmaps and every glyph in the
% PDF becomes a Type 3 font. That was measured: the first build produced a
% document with 10 Type 3 fonts and no Type 1. Latin Modern is Computer
% Modern redrawn with a full T1 repertoire and real .pfb outlines, so the
% same design comes out scalable.
\usepackage[T1]{fontenc}
\usepackage[utf8]{inputenc}
\usepackage{lmodern}
% microtype is not in the pinned TeX distribution. emergencystretch alone
% keeps the two-column measure from producing overfull lines.
\setlength{\emergencystretch}{3em}

\usepackage{amsmath}
\usepackage{amssymb}
\usepackage{amsthm}
\usepackage{booktabs}
\usepackage{array}
% enumitem is not in the pinned TeX distribution; the list spacing below
% does what \setlist would have done.
\usepackage{float}
\usepackage{xcolor}
\usepackage{listings}
\usepackage{tikz}
\usetikzlibrary{arrows.meta,positioning,fit,backgrounds,calc}
\usepackage[hidelinks,bookmarks=true]{hyperref}

% Numbered, definition-style box for a rule the reader is meant to apply.
\theoremstyle{definition}
\newtheorem{rulebox}{Rule}
\newtheorem{example}{Example}

% Unnumbered asides.
\theoremstyle{remark}
\newtheorem*{tip}{Tip}
\newtheorem*{pitfall}{Pitfall}

% Tighter lists, in place of enumitem's \setlist.
\setlength{\topsep}{2pt}
\setlength{\itemsep}{1pt}
\setlength{\parsep}{0pt}

\newcommand{\code}[1]{\texttt{#1}}
\newcommand{\term}[1]{\emph{#1}}

% Syntax highlighting. Four roles, distinguishable in colour and still
% distinguishable in greyscale, because an IEEE article gets printed: the
% keyword is the only bold role, the comment is the only italic one, and the
% two remaining roles differ in lightness as well as in hue.
\definecolor{kwblue}{RGB}{20,60,150}
\definecolor{cmtgreen}{RGB}{90,120,90}
\definecolor{strred}{RGB}{150,50,40}
\definecolor{numgray}{gray}{0.45}
\definecolor{framegray}{gray}{0.70}
\definecolor{bgshade}{gray}{0.97}
% A darker fill for figure cells. The listing background has to stay very
% light to keep code readable; a figure cell has to be clearly shaded or the
% distinction it draws is invisible in print.
\definecolor{cellshade}{gray}{0.90}

% The listing size is the document's choice, because it depends on the
% column: a two-column article sets `\listingsize` to 6pt and keeps its
% listings within 70 columns, a one-column document keeps the body size
% and includes source files at 80. Lines are never broken; a line that does not fit
% overflows the frame, which is visible, rather than wrapping, which is
% not.
\providecommand{\listingsize}{\scriptsize}

% `'` is deliberately not a string delimiter: the language uses it for loop
% labels, as Rust does, so treating it as a quote would swallow the rest of
% the line.
\lstdefinestyle{txhdl}{
  basicstyle=\ttfamily\listingsize,
  keywordstyle=\bfseries\color{kwblue},
  commentstyle=\itshape\color{cmtgreen},
  stringstyle=\color{strred},
  numberstyle=\tiny\color{numgray},
  morekeywords={mod,use,pub,const,type,struct,enum,trait,impl,fn,async,let,mut,
    match,if,else,loop,while,for,in,break,continue,return,as,await,
    module,bus,channel,transaction,protocol,pipeline,flow,seq,config,
    port,stage,pipe,instance,connect,bind,tag,domain,blackbox,foreign,
    property,role,signal,handler,true,false,
    sequence,interface,modport,implements,package,import,var,wire,func,proc,
    case,default,next,fork,branch,join,state,fsm},
  morecomment=[l]{//},
  morestring=[b]",
  showstringspaces=false,
  columns=fullflexible,
  keepspaces=true,
  breaklines=false,
  % Framed, so a listing is bounded rather than floating in the column.
  frame=single,
  rulecolor=\color{framegray},
  framesep=3pt,
  backgroundcolor=\color{bgshade},
  xleftmargin=3pt,
  xrightmargin=3pt,
  aboveskip=6pt,
  belowskip=6pt,
  % Every listing is numbered and captioned, so it can be referred to by
  % number. `caption` without `float` numbers the listing and leaves it
  % where it was written, which is what a listing in running prose needs.
  captionpos=b,
  abovecaptionskip=2pt,
  belowcaptionskip=6pt,
}

% Figures drawn as text. Framed the same way, and with the highlighting
% switched off, because a box-and-arrow diagram is not code and half its
% words turning blue reads as an accident. Setting a style adds keys rather
% than clearing the ones already set, so the three roles are neutralised by
% hand rather than by leaving them out.
\lstdefinestyle{ascii}{
  basicstyle=\ttfamily\listingsize,
  keywordstyle=\ttfamily,
  commentstyle=\ttfamily,
  stringstyle=\ttfamily,
  columns=fullflexible,
  keepspaces=true,
  frame=single,
  rulecolor=\color{framegray},
  framesep=3pt,
  backgroundcolor=\color{bgshade},
  xleftmargin=3pt,
  xrightmargin=3pt,
  aboveskip=6pt,
  belowskip=6pt,
}
\lstset{style=txhdl}

% House TikZ styles. `shorten` lives on the arrow styles rather than on each
% arrow, so every arrow in the document gets the gap, including ones added
% later. TikZ clips a path at a node's border, so without it an arrowhead
% butts against the box with no gap at all. Keep `node distance` well above
% twice the shorten, or the arrow shrinks to nothing.
\tikzset{
  box/.style={draw,rounded corners=2pt,align=center,inner sep=4pt,
              font=\footnotesize},
  boxf/.style={box,fill=bgshade},
  lbl/.style={font=\scriptsize\itshape,align=center},
  fl/.style={-{Stealth[length=4pt]},shorten >=2.5pt,shorten <=2.5pt},
  fld/.style={fl,dashed},
}



\InputIfFileExists{buildstamp.tex}{}{}
\input{timeline_counts}
\providecommand{\buildstamp}{Draft build (outside the Bazel flow).}

\title{The Dynamics of a Project}

\author{Filip Filmar and Dragiša Janković%
\thanks{Both authors are independent. Filip Filmar is the
corresponding author, at f@hdlfactory.com; Dragiša Janković is
reached at linkedin.com/in/dragisa-jankovic.}%
\thanks{This document is about how TxHDL was built rather than what
it is. What it is, is in the paper~\cite{paper}; the
cover~\cite{cover} lists every document. The chart and every number
here are measured from the repository's own history at build time, by
\code{//tools/timeline} and \code{//tools/gantt}, and not typed in.}%
\thanks{The authors used a large language model, Claude, as an
assistant in exploring the concepts and in writing and constructing
the documents and the programs; every commit of the repository says
so and carries the prompts that led to it, verbatim.}%
\thanks{\buildstamp}}

\markboth{The Dynamics of a Project}{Filmar and Janković: Dynamics}

\begin{document}
\maketitle

\begin{abstract}
A hardware description language, its runtime, its lowering to two
netlist languages, a RISC-V core that runs on an FPGA, an AXI bus
with a router, a small GPU, two compiler toolchains for the core, and
sixteen documents, were written in \ganttdays{} days of work spread
over four months. This is the shape of those \ganttdays{} days. It is
not a narrative of intentions: every bar in the chart is commits that
touched particular files on particular days, measured from the
repository at build time. Three things are visible in it that a plan
would not have predicted. The specification came four months before
the first line of the runtime, and nothing in between. The
implementation was not staged but simultaneous, because a standing
rule of the project requires every concept to enter the runtime, an
example and a document in the same change, so a feature is never a
bar in one band. And the parts that look like foundations, the bus
and the toolchains, are the youngest things in the tree.
\end{abstract}

\begin{IEEEkeywords}
Software process, project history, hardware description languages,
repository mining.
\end{IEEEkeywords}

\section{What Is Measured}
\label{sec:d-what}

A band of the chart is a set of paths, and a block is a day on which
a commit touched one of them, shaded by how many did. Paths rather
than commit messages, because a path is what a change actually moved
and a message is what it said it moved. The axis is the days the
project was worked on and not the calendar, for a reason
Section~\ref{sec:d-gap} gives.

What is not measured is the arrows. A repository records what changed
and never why it could not have changed sooner, so the dependencies
are stated in \code{//tools/timeline} by hand and are the one part of
the picture that is a claim rather than a measurement. They are the
claims the rest of the documents already make: that the lowering
needed values and units first, that the core needed the lowering,
that the board needed the core, that the router needed the link, and
that the C++ toolchain reused what the Rust one had built.

\section{The Four Month Gap}
\label{sec:d-gap}

The project has two lives. \ganttbefore{} commits over
\ganttbeforedays{} days in May are the specification: the language
argued out in prose, with no runtime under it. Then \ganttgapdays{}
days in which nothing was committed at all. Then \ganttafter{}
commits over \ganttafterdays{} days in September, which is
everything else in the tree. The gap is not stated here but found:
\code{//tools/gantt} takes the longest stretch between two working
days and counts what falls on either side of it.

That is why the chart's axis is working days rather than dates. On a
calendar the four idle months would be most of the picture and the
week that holds nine commits in ten would be a smudge a centimetre
wide. The gap is a fact about the project and it is said here in
words, where a reader can weigh it, rather than drawn as an expanse
of nothing.

The gap also explains why the specification band is so thin. It was
finished before the rest began, and it was touched three times
afterwards: twice to publish it as a hermetic article, and once for a
licence header. A document nobody edits is either finished or
forgotten, and which of the two it is cannot be measured from a
repository at all.

\begin{figure*}[t]
\centering
\input{gantt.tex}
\caption{The project by working day. A block is a day on which at
least one commit touched that band's files, shaded by how many did;
the count under each date is the commits of that day across the whole
tree. The arrows are what a band needed before it could begin. The
axis is working days and not the calendar, so the break marks where
it was cut and says what the cut stands for; every other step between
two columns is a night or a weekend. Both the cut's position and its
length are found from the history rather than chosen.}
\label{fig:d-gantt}
\end{figure*}

\section{Everything At Once}
\label{sec:d-together}

The striking thing in Fig.~\ref{fig:d-gantt} is how little of it
looks like a schedule. The runtime, the lowering, the examples and
the documents are all active on the same four days. A plan would have
staged them, and the plan would have been wrong.

The cause is a rule the project set itself before any of this was
written, stated in \code{AGENTS.md} and quoted in the
paper~\cite{paper}: a concept enters the runtime, a compiled example
and the documents in the same change, or the change is not finished.
The rule was adopted to stop documentation drifting from code. Its
effect on the shape of the work is that a feature is never a bar in
one band. Adding memories to the lowering, to take one day's work,
moved \code{lib/macros}, \code{lib/examples} and \code{docs} in the
same hour, and the chart can only show that as four bands lit at
once.

So the bands are not phases. They are layers, and the project moved
through all of them on every feature. Whether that is a good way to
work is not something a chart can answer; what the chart shows is
that the rule was kept.

The documents band is the widest and the darkest, with
\ganttdocs{} commits against \ganttcommits{} in the whole tree.
Some of that is the rule above. The rest is that this project
produces documents as its output and not as its paperwork: sixteen
of them, built from the tree at every commit, with every listing
included from a file and every figure drawn from a run.

\section{The Foundations Came Last}
\label{sec:d-last}

The bus, the router, the GPU and the two toolchains are all on the
last working day, and each of them is the kind of thing a plan puts
at the bottom.

That order is not an accident of enthusiasm. Each of them needed
something the project did not have until late. The AXI link needed
channels to be hardware, which is the \code{Buffer} part, which
needed the lowering to handle a send under a condition. The router
needed the link. The core moving onto AXI needed the router. The C++
toolchain needed the image tool and the linker script that the Rust
toolchain had made, and the second toolchain took a fraction of the
first's effort because of it.

The pattern is worth naming because it is the opposite of the usual
advice. The project did not build a bus and then a processor on it.
It built a processor on a bus of its own, ran it on an FPGA, and only
then wrote the standard bus, at which point the processor's own bus
was deleted. The core was the thing that proved the language, and a
standard bus is a convenience that a proof does not need.

\section{What A Repository Cannot Say}
\label{sec:d-cannot}

Three things are missing from the chart and no amount of mining would
recover them.

The first is time not spent committing. A day with seventy-one
commits and a day with two are two very different days, and the chart
says so, but neither says how long anyone sat at it.

The second is the work that failed. The tree keeps its dead ends on
purpose, as probes under
\code{experiments/rust\_embedding/}, and those are in the chart as
ordinary commits, indistinguishable from the work that succeeded. A
probe whose answer was no cost as much to write as one whose answer
was yes.

The third is order within a day. Almost every commit here landed on a
day with dozens of others, so the chart's smallest unit is far
coarser than the work's. An hour-by-hour axis was tried and
discarded: it is accurate and unreadable, which is the usual fate of
a chart that refuses to summarise.

\begin{thebibliography}{9}

\bibitem{paper}
F.~Filmar and D.~Janković, ``TxHDL: A hardware description language
that is a Rust library,'' project repository \code{HDL/txhdl},
\code{//docs:paper}, 2026.

\bibitem{cover}
F.~Filmar and D.~Janković, ``TxHDL documents,'' project repository
\code{HDL/txhdl}, \code{//docs:cover}, 2026.

\end{thebibliography}

\end{document}
